\documentclass[11pt,a4paper]{article}
\usepackage[margin=23mm,headheight=15pt,headsep=7mm,footskip=10mm]{geometry}
\usepackage[T1]{fontenc}
\usepackage[utf8]{inputenc}
\usepackage{amsmath,amsthm}
\usepackage{newtxtext,newtxmath}
\usepackage{booktabs,longtable,array,calc}
\usepackage{microtype,xurl,graphicx}
\usepackage[hidelinks]{hyperref}
\usepackage{fancyhdr}
\usepackage{titlesec}
\setlength{\parindent}{0pt}
\setlength{\parskip}{4pt}
\setlength{\emergencystretch}{2em}
\linespread{1.02}
\setcounter{secnumdepth}{0}
\setcounter{tocdepth}{2}
\newcommand{\tightlist}{\setlength{\itemsep}{0pt}\setlength{\parskip}{0pt}}
\providecommand{\pandocbounded}[1]{#1}
% Compatibility for uncaptioned Pandoc longtables with recent LaTeX.
\newcounter{none}
\pagestyle{fancy}
\fancyhf{}
\fancyhead[L]{\footnotesize Ptolemaic rigidity and stability}
\fancyhead[R]{\footnotesize Candidate v0.3.1}
\fancyfoot[C]{\thepage}
\renewcommand{\headrulewidth}{0.3pt}
\titleformat{\section}{\large\bfseries}{\thesection}{0pt}{}
\titleformat{\subsection}{\normalsize\bfseries}{\thesubsection}{0pt}{}
\titleformat{\subsubsection}{\normalsize\itshape}{\thesubsubsection}{0pt}{}
\hypersetup{pdftitle={Sharp local stability and critical-exponent
sensitivity of Ptolemaic split
metrics},pdfauthor={Anonymous},pdfsubject={Research proof supplement; unrefereed}}
\begin{document}
\begin{center}
{\LARGE\bfseries Sharp local stability and critical-exponent sensitivity
of Ptolemaic split metrics\par}
\vspace{6pt}
{\large All-cardinality first-order constants and the four-point
transition\par}
\vspace{8pt}
{\small Anonymous\quad |\quad Candidate v0.3.1\quad |\quad 27 September 2026\par}
\end{center}
\vspace{4pt}
\textbf{Assurance.} The new local theorems below are proved from
explicit matrix calculations and a finite collection of
rational-function certificates. They do not require the companion
all-size theorem. Calling the listed split points the complete equality
family for \(m\ge6\) uses the supplied v0.2 classification and its
stated dependencies. The new four-point classification imports the
four-point threshold, strict descent and Ercan's scalar estimate and
coefficient transport. No proof-assistant verification, authenticated
independent specialist review or exhaustive priority search is claimed.
The complete v0.2 archive is preserved unchanged.

\subsection{Abstract}\label{abstract}

We determine the sharp first-order stability constants at every binding
complete split metric of cardinality \(m\ge6\), and also at the
five-point split. Under mean-cross-distance normalisation, the
negative-type gap is comparable to the maximum distance perturbation
with exact lower and upper one-sided asymptotic constants. Eighteen
rational-function dual certificates prove the result in every larger
dimension; thirty-six finite certificates cover the small exceptions.
Their nonnegativity for the entire parameter range is certified by
polynomial coefficient expansions, not inferred from tests of finitely
many dimensions. We give explicit quadratic remainder bounds and prove
that every direction in the linearised active-constraint cone is
realised by a Ptolemaic curve. The resulting coefficient has three
stable parity-dependent formulas after cardinality eight. We then
calculate the first variation of the metric's own supremal negative-type
exponent, obtaining sharp directional sensitivity constants. The three-
and four-point equality families are also identified; at the four-point
star, unequal-arm perturbations produce a quadratic gap, so the linear
stability phenomenon does not extend to that cardinality. Finally, a
quotient argument identifies the full kernel at collision limits and an
explicit subcritical gap bound quantifies strictness for separated
smaller-cardinality metrics. No numerical value for the global
separated-class constant is asserted.

\subsection{1. Scope, notation and imported
results}\label{scope-notation-and-imported-results}

A symmetric array \(d\) on distinct labelled points is a Ptolemaic
metric if its distances are positive off the diagonal, satisfy the
triangle inequalities, and satisfy
\[d_{ij}d_{kl}\le d_{ik}d_{jl}+d_{il}d_{jk}.\] The metric
\(s=\operatorname{CS}(a,r)\) has parts \(A,B\) of sizes \(a,r\): its
distances are one within \(A\) and across the parts, and two within
\(B\). Put \(m=a+r\) and \[T=\frac{r(a+1)}{a(r-1)},\qquad p=\log_2 T.\]
The binding pairs considered here are \((a,r)=(k,k+1)\), \((k,k)\) and
\((k,k+2)\), with \(m\ge6\), together with \((2,3)\). Thus \(a\ge2\),
\(r\ge3\), and \(0<p<2\). For the first three families \(p\) is exactly
the universal threshold \(p_m\) in the companion paper; in the
five-point case this is Ercan's endpoint. The supplied v0.2 manuscript
classifies all equality metrics for \(m\ge6\) as similarities of these
splits. We do not re-prove that classification.

Let \(\Pi=I-\mathbf1\mathbf1^{\mathsf T}/m\) and define
\[\Delta_p(d)=\min_{z\perp\mathbf1,\ \|z\|_2=1}-z^{\mathsf T}d^p z.\]
Powers are entrywise. This is a restricted eigenvalue; the automatic
centring zero on the full space is not included. At the split, the
restriction of \(-s^p\) has eigenvalues \(0\), \(1\) with multiplicity
\(a-1\), and \(T\) with multiplicity \(r-1\). The unit null direction is
proportional to \((r\mathbf1_A,-a\mathbf1_B)\), with squared normalising
factor \(arm\).

Normalise the mean of the \(ar\) cross-distances to one. Write
\(e=d-s\), with coordinates \(y_{ij}\) within \(A\), \(x_{ib}\) across,
and \(t_{bc}\) within \(B\), and set \(h=\|e\|_\infty\). The
normalisation is \(\sum x_{ib}=0\). Throughout, a local statement is at
a fixed labelled split; constants may depend on \((a,r)\).

\textbf{Attribution.} The transport of a zero-sum coefficient vector
through metric inversion is Ercan's Theorem 8.1, and the sign-changing
choice of vertex is his Theorem 8.2 {[}E{]}. The equality-potential
observation in v0.2 is used below, with these attributions retained. The
threshold and Lorentzian correspondence originate in {[}B{]}; strict
negative-type descent is {[}LW{]}. Matrix perturbation is standard
{[}K{]}. The new work in this supplement concerns the complete
tangent-cone calculation, the all-size sharp constants and the
critical-exponent expansion. The short small-cardinality consequences
are not presented as independently verified historical novelties.

\subsection{2. Main local theorem}\label{main-local-theorem}

Define the sum of all linearised active slacks by
\[\ell(e)=-r(r-1)\sum_{i<j\in A}y_{ij}
-\frac{a(a+1)}2\sum_{b<c\in B}t_{bc}.\] Set
\[\gamma=\frac{2p}{ma(r-1)},\quad
U_{a,r}=\frac{pr(3a-1)}{2m},\quad A_{a,r}=\gamma c_{a,r},\] where the
rational constant \(c_{a,r}\) is

{\def\LTcaptype{none} % do not increment counter
\begin{longtable}[]{@{}lll@{}}
\toprule\noalign{}
Pair or family & Range & \(c_{a,r}\) \\
\midrule\noalign{}
\endhead
\bottomrule\noalign{}
\endlastfoot
\((2,3)\) & five-point local result & \(6/5\) \\
\((2,4)\) & six points & \(3/2\) \\
\((3,3)\) & six points & \(9/5\) \\
\((3,4)\) & seven points & \(3\) \\
\((3,5)\) & eight points & \(60/11\) \\
\((4,4)\) & eight points & \(5\) \\
\((k,k)\) & \(k\ge5\) & \(k^2/2\) \\
\((k,k+1)\) & \(k\ge4\) & \(k(k+1)/2\) \\
\((k,k+2)\) & \(k\ge4\) & \(k(k+1)/2\) \\
\end{longtable}
}

\textbf{Theorem 2.1 (sharp first-order stability in all these
dimensions).} Put \[b_{a,r}=p(m-1)+18p^2(m-1)^2+\frac{3pr(a-1)}m,
\qquad \rho_{a,r}=\min\left\{\frac18,\frac1{12p(m-1)}\right\}.\] Every
Ptolemaic \(d\) in the preceding normalisation with \(h\le\rho_{a,r}\)
satisfies \[\boxed{A_{a,r}h-b_{a,r}h^2\le\Delta_p(d)
\le U_{a,r}h+p(m-1)h^2.}\] Both linear coefficients are sharp:
\[\liminf_{\substack{d\to s,\ d\ne s\\d\ {\rm Ptolemaic}}}\frac{\Delta_p(d)}h=A_{a,r},
\qquad
\limsup_{\substack{d\to s,\ d\ne s\\d\ {\rm Ptolemaic}}}\frac{\Delta_p(d)}h=U_{a,r}.\]
In particular \(\Delta_p(d)\ge A_{a,r}h/2\) if also
\(h\le A_{a,r}/(2b_{a,r})\). The least eigendirection is simple on the
displayed neighbourhood; with consistent sign its distance from the
split null vector is at most \(6p(m-1)h\).

The general quadratic constants are deliberately conservative. They do
not replace the substantially better seven-point remainder and radius
already proved in v0.2. The advance is exact first-order control in
every cardinality and every equality type, rather than a numerical
improvement of that seven-point radius.

For \(m\ge9\), with mean cross-distance one and \(h=\|d-s\|_\infty\)
still fixed, the lower coefficient simplifies to \[\boxed{
A_{a,r}=\begin{cases}
\displaystyle p_m\frac{m+1}{m(m-1)},&m\text{ odd},\\[1mm]
\displaystyle\frac{p_m}{m-2},&m\text{ even},\ a=r,\\[1mm]
\displaystyle\frac{p_m}{m},&m\text{ even},\ r=a+2.
\end{cases}}\] Thus the two even equality types have different local
condition constants. The metric normalisation and the norm used here are
part of the statement; these constants are not invariant under changing
that convention.

\subsection{3. The exact tangent cone}\label{the-exact-tangent-cone}

At \(s\), the only active triangle and Ptolemy inequalities are
\[u_{i;bc}=x_{ib}+x_{ic}-t_{bc}\ge0,\]
\[v_{ij;bc}=x_{ib}+x_{ic}+x_{jb}+x_{jc}-2y_{ij}-t_{bc}\ge0,\] for
\(i,j\in A\), \(b,c\in B\), with unordered distinct pairs. There are
\(a\binom r2\) triangle rows and \(\binom a2\binom r2\) Ptolemy rows.
Denote their integer matrix by \(D\). If \(w\) is the cross-edge
indicator, direct counting gives
\[D^{\mathsf T}\mathbf1=a(r-1)w-r(r-1)\mathbf1_{AA}
-\frac{a(a+1)}2\mathbf1_{BB}.\] Hence \(\ell(e)=\mathbf1^{\mathsf T}De\)
when \(w^{\mathsf T}e=0\).

The exact metric inequalities give \(u\ge0\) and \[v_{ij;bc}\ge-3h^2,\]
because the omitted quadratic term is
\[x_{ib}x_{jc}+x_{ic}x_{jb}-y_{ij}t_{bc}.\] All inactive inequalities
have strictly positive slack at \(s\).

\textbf{Lemma 3.1 (every cone direction is geometrically attainable).}
Let \(De\ge0\), \(w^{\mathsf T}e=0\), and \(\|e\|_\infty\le1\). Let
\(e^{\rm in}\) equal \(-1\) on within-part edges and zero on
cross-edges. Then \[d_t=s+te+4t^2e^{\rm in},\qquad 0<t\le1/20,\] is a
positive Ptolemaic metric with mean cross-distance one.

\emph{Proof.} On triangle rows \(De^{\rm in}=1\) and on Ptolemy rows
\(De^{\rm in}=3\). Consequently the active triangle slack is at least
\(4t^2\), and the linear part of each active Ptolemy slack is at least
\(12t^2\). The absolute quadratic contribution is at most
\(3(t+4t^2)^2\le(108/25)t^2<12t^2\). Every entry changes by at most
\(t+4t^2\le3/50\). An inactive triangle slack changes by at most
\(9/50<1\). An inactive Ptolemy slack, a signed sum of three products,
changes by at most \(3(4\cdot3/50+(3/50)^2)<1\). The original inactive
slacks are at least one, so all remain positive. Positivity of distances
and the normalisation are immediate. Conversely, the derivative of any
differentiable feasible curve satisfies the linear inequalities. Thus
this is the exact tangent cone, not merely a relaxation. \(\square\)

\subsection{4. The sharp linear programme and its universal
certificate}\label{the-sharp-linear-programme-and-its-universal-certificate}

\textbf{Theorem 4.1 (exact cone norm).} For every binding pair in
Section 2, \[De\ge0,\quad w^{\mathsf T}e=0
\quad\Longrightarrow\quad
c_{a,r}\|e\|_\infty\le\ell(e)
\le\frac{ar(r-1)(3a-1)}4\|e\|_\infty.\] Both constants are attained by
cone directions.

The upper bound follows from the explicit coefficients of \(\ell\), and
\(e^{\rm in}\) attains it. The lower bound has a finite rational
certificate for every edge \(f\) and sign \(\sigma\in\{-1,1\}\):
\[\boxed{D^{\mathsf T}\eta+\mu w
=D^{\mathsf T}\mathbf1-c_{a,r}\sigma\epsilon_f,
\qquad\eta\ge0.} \tag{C}\] The supplied certificates additionally
satisfy \(\eta_j\le2\) on every Ptolemy row. Consequently the same
identity for a nonlinear feasible metric gives
\[\boxed{\ell(e)\ge c_{a,r}h
-6\binom a2\binom r2h^2.} \tag{4.1}\] Indeed choose \(\sigma e_f=h\).
Triangle slacks are nonnegative, Ptolemy slacks are at least \(-3h^2\),
and the sum of their certificate weights is at most
\(2\binom a2\binom r2\).

\subsubsection{4.1 Exact description of the certificate
families}\label{exact-description-of-the-certificate-families}

The certificate is invariant under the stabiliser of the selected edge.
For an anchor edge distinguish its two anchors; for a cross edge
distinguish its anchor and pole; for a pole edge distinguish its two
poles. A triangle row is classified by the number of distinguished
anchors and poles it contains. A Ptolemy row is classified in the same
way. Denote these row orbits by \((u,i,j)\) and \((v,i,j)\).

The data in \texttt{certificates/symbolic\_duals.json} give eighteen
explicit rational-function vectors \(\eta(k)\) and scalars \(\mu(k)\):
three families \((a,r)=(k,k+d)\), \(d=0,1,2\), three edge types and two
signs. The domains are \(k\ge5\) for \(d=0\) and \(k\ge4\) for
\(d=1,2\). \texttt{CERTIFICATE\_TABLES.md} prints all these expressions.
No linear-programming solver is required to verify them.

\textbf{Indispensable proof supplement.} The complete tables, both JSON
certificate collections, corrected verifiers, pinned dependencies and
checksum manifest accompany this version in the archived research
package {[}SUPP{]}. They are part of the proof of Theorem 4.1, not
optional numerical illustrations. The eighteen symbolic certificates
establish universal parameter ranges; sampled dimensions alone do not.

For clarity, one extremal certificate, the negative pole-edge
certificate in the balanced family \((k,k)\), is

{\def\LTcaptype{none} % do not increment counter
\begin{longtable}[]{@{}ll@{}}
\toprule\noalign{}
Row orbit & Weight \\
\midrule\noalign{}
\endhead
\bottomrule\noalign{}
\endlastfoot
\((u,0,0)\) & \(k/(k-3)\) \\
\((u,0,1)\) & \(0\) \\
\((u,0,2)\) & \(0\) \\
\((v,0,0)\) & \((k^2-4k-3)/[(k-3)(k-1)]\) \\
\((v,0,1)\) & \((k+1)/(k-1)\) \\
\((v,0,2)\) & \(1/(k-1)\) \\
\end{longtable}
}

Here \(\mu=1\) and \(c=k^2/2\). At \(k=5+x\), every weight is
nonnegative for \(x\ge0\); for example \(k^2-4k-3=x^2+6x+2\). All
Ptolemy weights are at most two.

For each of the eighteen vectors, the symbolic verifier constructs the
orbit matrix directly by counting. If \(s_A,s_B\) distinguished vertices
occur in the two parts, an orbit has multiplicity
\[\binom{s_A}{i}\binom{a-s_A}{n_A-i}
\binom{s_B}{j}\binom{r-s_B}{2-j},\] where \(n_A=1\) for a triangle row
and \(2\) for a Ptolemy row. The coefficient of an edge orbit in
\(D^{\mathsf T}\eta\) is its total coefficient in a representative row,
times this multiplicity, divided by the size of the edge orbit. Symmetry
proves this counting formula. Substitution verifies (C) as a
rational-function identity. Nonnegativity of \(\eta\) and of \(2-\eta\)
for Ptolemy rows is checked by substituting \(k=k_0+x\): each numerator
and denominator has nonnegative rational coefficients and the
denominator has positive constant term. These are certificates on the
entire half-line \(x\ge0\), not interpolation from sampled dimensions.

The six small pairs have thirty-six rational certificate vectors in
\texttt{small\_duals.json}, again covering the three edge types and both
signs, all satisfying the same Ptolemy weight bound. A separate
standard-library implementation enumerates the original labelled rows
and verifies the identities without using symbolic orbit counting. It
also evaluates the universal certificates at selected dimensions as a
separate implementation check. The symbolic positivity certificates, not
those finite evaluations, justify arbitrary cardinality.

\subsubsection{4.2 Primal sharpness
certificates}\label{primal-sharpness-certificates}

Distinguish two poles. Give every anchor the same cross perturbation
\(u\) towards these two poles and \(v=-2u/(r-2)\) towards the others.
Give all anchor edges the same perturbation \(Y\), the distinguished
pole edge perturbation \(-1\), each pole edge with exactly one
distinguished endpoint perturbation \(t_1\), and every other pole edge
perturbation \(t_2\).

For the six small pairs use
\[u=-\frac{r-2}{3r-4},\quad v=\frac2{3r-4},\quad
Y=-\frac{r-4}{2(3r-4)},\quad t_1=-\frac{r-4}{3r-4},\quad t_2=\frac4{3r-4}.\]
These entries have maximum absolute value one, obey \(De\ge0\) and the
normalisation, and give
\[\ell(e)=\frac{ar\{(a-1)(r-2)(r-3)+4\}}{4(3r-4)}=c_{a,r}.\] There is no
\(t_2\) edge when \(r=3\); the unused value causes no issue.

For the balanced family \(a=r=k\ge5\) use
\[u=-\frac{r-2}{4(r-1)},\quad v=Y=\frac1{2(r-1)},\quad
 t_1=-\frac{r-2}{2(r-1)},\quad t_2=\frac1{r-1}.\] The six row-orbit
slacks, in increasing numbers of distinguished poles, are
\(0,r/[4(r-1)],r/[2(r-1)]\) for triangles and \(0,0,0\) for Ptolemy
rows. Every entry has absolute value at most one, and direct
substitution gives \(\ell=ar/2=c_{a,r}\). In the remaining two infinite
families, simply shorten one pole edge by one and leave every other
perturbation zero; this gives \(\ell=a(a+1)/2=c_{a,r}\). Thus the lower
bound is attained. Lemma 3.1 integrates these directions into true
metrics, proving that the linear-programming constants are the geometric
first-order constants. \(\square\)

\subsection{5. From the cone to the spectral
gap}\label{from-the-cone-to-the-spectral-gap}

The exceptional constants reflect a change of the optimal face, not a
different choice of normalisation. In the six small pairs, shortening
the distinguished pole edge is accompanied by compensating cross-edge
and anchor-edge changes. Their cost is smaller than that of shortening
that edge alone. For example, at \((3,5)\) this gives \(60/11<6\). In
the stable unbalanced families the single-edge direction is already
optimal. In the balanced family the optimum continues to distribute the
perturbation, with all Ptolemy rows tight in the displayed stable
direction. Its dual certificate requires \(k\ge5\): the numerator
\(k^2-4k-3\) is negative at \(k=4\). These explicit primal and dual
facts explain why the small regimes require separate certificates; they
are not a claim of uniqueness of the extremising direction.

The null eigenvalue at \(s\) is simple on \(\mathbf1^\perp\), and the
other eigenvalues are at least one. Put \(H=-\Pi d^p\Pi\) on that space,
and expand \(H=H_0+H_1+H_2\). Because \(0<p<2\), \(h\le1/8\), and the
reference distances are one or two, entrywise Taylor expansion gives
\[\|H_1\|\le2p(m-1)h,\qquad
\|H_2\|\le p(m-1)h^2,\qquad
\|H-H_0\|\le3p(m-1)h=:\varepsilon.\] The second derivative bound used
here is \(|p(p-1)x^{p-2}|/2\le p\) on \(x\ge7/8\) when \(0<p<2\) and
\(x\) lies within \(1/8\) of one or two. This loose bound also covers
\(p=1\).

Direct expansion against the unit null vector gives
\[z_*^{\mathsf T}H_1z_*
=-\frac{2p}{m}\left[\frac ra\sum y_{ij}
+\frac{aT}{2r}\sum t_{bc}\right]=\gamma\ell(e).\] The cross term
vanishes because the sum of cross perturbations is zero. This does not
mean that individual cross edges are irrelevant: the active constraints
couple them to within-part changes.

If \(\varepsilon\le1/4\), the other eigenvalues remain at least \(3/4\),
and the least eigenvalue remains in \([-1/4,1/4]\). Write the matrix in
the decomposition \(\mathbb Rz_*\oplus z_*^\perp\). The scalar Schur
equation gives \[0\le z_*^{\mathsf T}Hz_*-\lambda_{\min}(H)
\le\frac{\varepsilon^2}{1-2\varepsilon}\le2\varepsilon^2.\] For the
corresponding signed unit eigenvector, the block equation bounds the
tangent of its angle from \(z_*\) by \(\varepsilon/(1-2\varepsilon)\).
The distance between the unit vectors is no larger than that tangent,
hence at most \(2\varepsilon\). This is the elementary simple-eigenvalue
estimate underlying standard perturbation theory {[}K{]}. Therefore
\[\gamma\ell(e)-[p(m-1)+18p^2(m-1)^2]h^2
\le\Delta_p(d)\le\gamma\ell(e)+p(m-1)h^2.\] Insert (4.1) and the upper
cone bound. Since \[6\gamma\binom a2\binom r2=\frac{3pr(a-1)}m,\]
Theorem 2.1 follows. The extremal curves of Section 4 have
\(h(t)=t+O(t^2)\) and \(\ell(e(t))=c_{a,r}t+O(t^2)\) or its upper
analogue, proving sharpness.

This proof is local and self-contained at the specified split metrics.
It gives local negative type without assuming the global companion
theorem. The companion and v0.2 remain premises only when one concludes
that every possible global equality metric has one of these local
models.

\subsection{6. Sensitivity of the metric's own critical
exponent}\label{sensitivity-of-the-metrics-own-critical-exponent}

Let \[\wp(d)=\sup\{q>0:d\text{ has }q\text{-negative type}\}.\] This is
the exponent of the particular metric, not the universal minimum over
all metrics of its cardinality.

\textbf{Relation to earlier exponent and gap results.} Sánchez {[}S,
Theorem 2.3 and Corollary 2.4{]} characterises the finite-metric
endpoint using the distance matrix and its inverse on the constant
vector. Thus detecting the critical exponent through loss of matrix
strictness is established machinery. His complete-bipartite path metrics
are not the complete split metrics here. Li--Weston {[}LW, Theorem
3.3{]} give quantitative extensions of the admissible exponent from a
positive negative-type gap, as well as strict descent and endpoint
non-strictness. The contribution here is the exact Ptolemaic tangent
cone, its attained extremal directional constants, and their
all-cardinality formulas, not a new general matrix characterisation.

Gap normalisations must also be distinguished. Our \(\Delta_p\)
minimises \(-z^{\mathsf T}d^p z\) over zero-sum vectors with
\(\|z\|_2=1\). The normalised-simplex gap in {[}LW{]} uses positive and
negative weights separately summing to one (equivalently \(\|z\|_1=2\)),
with its associated quadratic-form factor. A rooted Gramian has yet
another coordinate convention; Robertson {[}R, Sections 1--2{]}
discusses its spectral relation to negative-type gaps for ultrametric
spaces. These quantities express related strictness properties but their
numerical constants cannot be substituted without a normalisation
conversion.

\textbf{Theorem 6.1 (critical-exponent expansion).} At every split
covered by Theorem 2.1, for Ptolemaic \(d\to s\) in mean-cross-distance
normalisation,
\[\boxed{\wp(d)-p=\frac{m}{r(a+1)\log2}\,\Delta_p(d)+O(h^2)
=\frac{2p\ell(e)}{ar(r-1)(a+1)\log2}+O(h^2).}\] The constants hidden in
\(O(h^2)\) depend on the fixed split. Consequently the sharp lower and
upper limits of \((\wp(d)-p)/h\) are
\[\boxed{\frac{2p c_{a,r}}{ar(r-1)(a+1)\log2},
\qquad\frac{p(3a-1)}{2(a+1)\log2}.}\]

\emph{Proof.} In positive distance coordinates, \(-\Pi d^q\Pi\) is real
analytic jointly in \(d,q\). Its simple eigenvalue through zero at
\((s,p)\) has exponent derivative \[\partial_q\lambda(s,p)
=-\frac{a(r-1)T}{m}\log2
=-\frac{r(a+1)}m\log2=:-\Lambda<0.\] The implicit-function theorem
produces a unique nearby zero \(q_c(d)\), with other restricted
eigenvalues positive. The least eigenvalue changes from positive to
negative as \(q\) increases through this root. Strict descent of
negative type {[}LW{]} prevents another, larger negative-type interval:
negative type at any larger exponent would imply strict negative type at
an exponent immediately above the displayed root, a contradiction. Hence
\(\wp(d)=q_c(d)\). Expansion at \(q=p\) gives
\[0=\Delta_p(d)-\Lambda(q_c-p)+O(h|q_c-p|+(q_c-p)^2),\] and
\(q_c-p=O(h)\) gives the first formula. Theorem 2.1 and its sharp curves
give both limits. \(\square\)

For seven points these limits are
\[\frac{p_7}{24\log2}\quad\text{and}\quad\frac{p_7}{\log2}.\] This
identifies how much extra snowflaking power a near-extremal metric
actually permits. It is not an algorithmic complexity statement. The
expansion concerns a fixed cardinality and normalisation; no uniform
remainder over all \(m\) is claimed.

For large \(m\), \(p_m\sim4/(m\log2)\). The sharp lower gap coefficient
is asymptotic to \(4/(m^2\log2)\), while the upper coefficient tends to
\(3/(2\log2)\). Their ratio grows as \(3m^2/8\). The sharp lower
critical-exponent coefficient is asymptotic to \(16/[m^3(\log2)^2]\),
whereas its upper coefficient is asymptotic to \(6/[m(\log2)^2]\). These
follow by substitution, not by extrapolating numerical data.

\subsection{7. Small-cardinality equality and the change of stability
order}\label{small-cardinality-equality-and-the-change-of-stability-order}

\textbf{Proposition 7.1 (three and four points).} At \(P(3)=2\), the
equality metrics are precisely the three distinct collinear points,
equivalently the positive triples whose longest distance is the sum of
the other two. At \(P(4)=\log_2 3\), the equality metrics are precisely
the similarities of \(\operatorname{CS}(1,3)\) and
\(\operatorname{CS}(2,2)\).

\emph{Proof.} The three-point assertion is the Gram determinant
criterion for a Euclidean triangle. For four points put \(p=\log_2 3\).
The negative-type endpoint is {[}B, Theorem 7.16{]}; every three-point
metric has strict negative type at this \(p<2\) by {[}LW{]}. Thus an
equality vector has full support.

If its sign pattern is \(1+3\), normalise the positive coefficient to
one and write the negative magnitudes as \(c_i>0\), \(\sum c_i=1\). Let
the three root lengths be \(\ell_i\). Ercan's scalar inequality {[}E,
Lemma 2.1{]} with \(2^p-2=1\) gives
\[(\ell_i+\ell_j)^p\le\ell_i^p+\ell_j^p+(\ell_i\ell_j)^{p/2},\] with
equality only when \(\ell_i=\ell_j\). The gap is a sum of the
triangle-power slacks, these scalar slacks and
\[\frac12\sum_{i<j}(c_i\ell_i^{p/2}-c_j\ell_j^{p/2})^2.\] All weights
are positive. Equality therefore forces equal root lengths, leaf
distances equal to their sums, and equal negative weights: the equal-arm
star \(\operatorname{CS}(1,3)\).

For a \(2+2\) vector, negative type implies \(d^p z=\lambda\mathbf1\).
The potential \(\lambda\) is nonzero: otherwise Ercan's transport at any
vertex would produce a nonzero equality vector supported on only three
points of an inverted Ptolemaic metric, contradicting their strictness.
Choose a vertex whose coefficient has the sign of \(\lambda\). Its
transported sign flips, giving a \(1+3\) pattern after a global sign
change. The inverted metric is the equal-arm star. Inverting back at the
relevant independent vertex gives \(\tfrac12\operatorname{CS}(2,2)\),
hence the second similarity type. Conversely both displayed metrics
attain equality by direct substitution. The transport and the
sign-selection mechanism are {[}E, Theorems 8.1--8.2{]}; only their
equality specialisation is being used. \(\square\)

This elementary consequence does not settle five-point equality. In
particular we do not replace Ercan's system of five-point equality
conditions by an unsupported classification. The local theorem above is
valid at \(\operatorname{CS}(2,3)\), regardless of whether another
five-point equality family exists.

\textbf{Proposition 7.2 (a quadratic four-point direction).} Put
\(p=\log_2 3\). Take the four-point star with root lengths \(1+t,1-t,1\)
and leaf distances their pairwise sums. Its mean root length is one and
\(h=|t|\). Then \[\boxed{\Delta_p(d_t)=\frac{p(2p-3)}8t^2+O(t^4).}\] The
coefficient is positive because \(3>2^{3/2}\). Therefore no positive
local linear lower bound \(\Delta_p(d)\ge c h\) can hold near the
four-point star.

\emph{Proof.} The equality vector is \((3,-1,-1,-1)/\sqrt{12}\); the
other restricted eigenvalues are both three. The first scalar variation
is zero. The direct quadratic Rayleigh term is \(3p(p-1)t^2/8\), and the
second-order eigenvector correction is \(p^2t^2/8\). Their difference is
the coefficient displayed. Swapping the first two leaves changes \(t\)
to \(-t\), so the simple eigenvalue is an even analytic function and the
next term is \(O(t^4)\). The symbolic verifier reconstructs these terms
from the six raw distances, not from the asserted coefficient.
\(\square\)

This is an obstruction to linear stability at four points, not a proof
of a general quadratic lower stability theorem there. The different
local order is a genuine distinction between that endpoint and the
higher-cardinality split metrics.

\subsection{8. Collision limits and an explicit smaller-cardinality
gap}\label{collision-limits-and-an-explicit-smaller-cardinality-gap}

The classification can be stated cleanly on the compactification,
without mislabelling collisions as equality metrics on distinct points.

\textbf{Proposition 8.1 (kernel at a collision limit).} Let a
diameter-one Ptolemaic pseudometric on \(m\ge6\) labelled points have
\(N<m\) equivalence classes under zero distance. At exponent \(p_m\),
the null space of its restricted negative-type form consists exactly of
vectors whose coefficient sum is zero separately in each equivalence
class. Its dimension is \(m-N\).

\emph{Proof.} The quotient is an \(N\)-point Ptolemaic metric. Since
\(p_m<p_N\) for \(N\ge3\), the supplied threshold results and strict
descent give strict \(p_m\)-negative type on the quotient. The \(N=2\)
case is immediate. For any vector, the original quadratic form equals
the quotient quadratic form applied to the class sums. On the zero-sum
subspace it vanishes exactly when every class sum vanishes. \(\square\)

Thus, conditional on the global threshold and classification, the zero
set on the diameter-one compactification comprises exactly the collision
strata and the finitely many non-colliding split copies. The collision
kernel has dimension \(m-N\); a genuine extremiser has one nontrivial
null direction, of full support.

A useful quantitative input for a future explicit global constant is
available without a compactness argument.

\textbf{Proposition 8.2 (explicit strict-descent gap).} Suppose an
\(N\)-point metric \(d\), \(N\ge2\), has \(q\)-negative type and minimum
distance at least \(\eta>0\). If \(0<p<q\) and \(\alpha=p/q\), then
\[\boxed{\Delta_p(d)\ge
\frac{\eta^p}{2\Gamma(1-\alpha)\,[\log(2(N-1))]^\alpha}.}\]

\emph{Proof.} Realise \(d^q\) as squared Euclidean distances. The
Gaussian kernel \(K_t=(e^{-td_{ij}^q})\) is positive semidefinite: for a
Euclidean realisation \(d_{ij}^q=\|x_i-x_j\|^2\), expand
\(e^{2t\langle x_i,x_j\rangle}\) as a nonnegative sum of positive
semidefinite Gram powers, then multiply by the positive diagonal factors
\(e^{-t\|x_i\|^2}\). For \(t\ge t_0=\log(2(N-1))/\eta^q\), every
off-diagonal entry is at most \(1/[2(N-1)]\), so \(K_t\succeq I/2\) by
diagonal dominance. For a unit zero-sum vector, \[-z^{\mathsf T}d^p z
=\frac{\alpha}{\Gamma(1-\alpha)}
\int_0^\infty z^{\mathsf T}K_tz\,t^{-1-\alpha}\,dt
\ge\frac{t_0^{-\alpha}}{2\Gamma(1-\alpha)}.\] This is the asserted
bound. The Gaussian positivity and integral representation are the
standard strict-descent mechanism; the explicit tail estimate is the
calculation used here. \(\square\)

For example, proper six-point submetrics at exponent \(p_7\) admit this
bound with \(q=1\), under the six-point threshold premise. This removes
one source of non-effective constants in a quantitative inverse proof.
It does not by itself yield a numerical \(C_\eta\) for the full
seven-point problem. An explicit global constant and uniform
collision-versus-rigidity estimates remain open tasks in this
supplement.

\subsection{9. Verification, novelty boundaries and next
obligations}\label{verification-novelty-boundaries-and-next-obligations}

The supplied v0.2 archive replayed in full mode with the nested
companion exact replay. Its exact, symbolic and LP reports matched the
supplied files; the search trial records matched, while the search
summary differed in runtime metadata. The unchanged replay receipt
records the difference. We did not authenticate the reviewers or treat
their reports as external specialist endorsements.

The new verification has three distinct roles:

\begin{enumerate}
\def\labelenumi{\arabic{enumi}.}
\tightlist
\item
  The symbolic script verifies all eighteen universal rational-function
  certificates by direct orbit counts. Positive-coefficient expansions
  prove the required inequalities for every real parameter \(k\ge k_0\),
  hence every integer cardinality in the stated families. It also
  re-derives the first-variation and four-point quadratic coefficients.
\item
  A separate standard-library rational script enumerates original
  labelled constraints, checks the thirty-six small certificates and
  seventy-two sampled evaluations of the general certificates, verifies
  all sharp primal directions and their exact Ptolemaic integrations,
  and checks the small-star examples. The sampled dimensions are
  implementation checks only.
\item
  Numerical perturbation and critical-exponent calculations corroborate
  the analytic formulas. The seeded convex-mixture trials combine only
  the displayed lower extremiser with the inward upper extremiser for
  each split; they do not sample the full feasible cone. They neither
  establish universal bounds nor substitute for an equality
  classification. High-precision records are retained separately.
\end{enumerate}

In version 0.3.1 all proof-critical Python assertions have been replaced
by explicit exception checks. Direct and replay entry points are tested
under normal Python, \texttt{-O}, \texttt{-OO}, and both optimisation
environment settings. Disposable negative controls alter a multiplier or
the exact identity and regenerate the manifest, so checksum failure
cannot mask the mathematical test. The retained mode-control report
states the executed scope. These software checks do not formally verify
the analytic proofs.

The dependency split is: Sections 3--5 use direct geometry, exact
certificates and the written spectral estimate; Section 6 adds the
classical endpoint/strict-descent results; Section 7 imports the public
small-cardinality results and Ercan's transport; Section 8 imports the
companion's threshold premises for the quotient interpretation. Only the
claim that the local models exhaust all global equality cases
additionally uses the retained v0.2 classification.

The certificate tables are part of the proof evidence, not merely
exploratory solver output. No optimisation solver is needed by either
exact verifier. Numerical linear programming was used during discovery,
and its exploratory scripts and logs are retained under
\texttt{provenance/discovery/} rather than described as proof
certificates.

The following remain unproved here: a classification of five-point
equality, a full quadratic stability theorem at four points, a numerical
global separated-class constant, or a dimension-uniform sharp
second-order remainder. The critical-exponent theorem has a proved
fixed-cardinality expansion, not a numerical claim about all
neighbourhood sizes. The local all-size result does not require the
unrefereed companion; the global equality interpretation does.
Historical priority has not been established for any new lemma.

\subsection{References}\label{references}

\subsubsection{Source manuscripts and immediate
predecessors}\label{source-manuscripts-and-immediate-predecessors}

{[}B{]} Baker, M., Huh, J., Kummer, M., \& Lorscheid, O. (2026).
\emph{Lorentzian polynomials and matroids over triangular hyperfields.
Part 2: Analytic aspects} (Version 2) {[}Preprint{]}. arXiv.
\url{https://arxiv.org/abs/2607.15375v2}

{[}E{]} Ercan, E. (2026). \emph{Ptolemaic negative type and the values
q(6) and q(7)} (Version 1) {[}Preprint{]}. arXiv.
\url{https://arxiv.org/abs/2608.20606v1}

Companion manuscript. (2026). \emph{Sharp negative type of finite
Ptolemaic metrics} (Unrefereed version 1.0). Supplied Evidence Press
research bundle, retained inside the immutable v0.2 input.

Supplied research manuscript. (2026). \emph{Rigidity at the sharp
Ptolemaic negative-type threshold} (Research version 0.2). Supplied 27
September 2026; retained unchanged in the input archive under
\texttt{provenance/}.

\subsubsection{Matrix perturbation and strict
descent}\label{matrix-perturbation-and-strict-descent}

{[}K{]} Kato, T. (1995). \emph{Perturbation theory for linear operators}
(Reprint of the 1980 edition). Springer.
\url{https://doi.org/10.1007/978-3-642-66282-9}

{[}LW{]} Li, H., \& Weston, A. (2010). Strict p-negative type of a
metric space. \emph{Positivity, 14}(3), 529--545.
\url{https://doi.org/10.1007/s11117-009-0035-2}. Consulted preprint:
\url{https://arxiv.org/abs/0901.0695}

{[}S{]} Sánchez, S. (2011). \emph{On the supremal p-negative type of a
finite metric space} {[}Preprint{]}.
\url{https://arxiv.org/abs/1108.0451}

{[}R{]} Robertson, G. (2025). \emph{Spectral properties of the Gramian
of finite ultrametric spaces} {[}Preprint{]}.
\url{https://arxiv.org/abs/2504.14840}

{[}SUPP{]} Anonymous (2026). \emph{Ptolemaic split stability: versioned
proof supplement}, v0.3.1-candidate. Complete tables, certificates, code
and retained predecessor. Repository:
\url{https://github.com/ipitchford/ptolemaic-split-stability} . Version
DOI: \url{https://doi.org/10.5281/zenodo.23000081}
\end{document}
